\reportheading{Appendix E. Manufacturer air-distribution assessment}
\subhead{E.1 Application to the specified diffuser}
The main design uses eighty Krueger 1400/51400 ceiling diffusers. Appendix A reports catalogue throw divided by the corresponding plan length. The manufacturer assessment examines the Type 3 circular-pattern curves in Krueger Air Distribution Engineering [4]. Each curve relates throw ratio to predicted ADPI for a specified sensible load.

\tabletitle{Table E1. Manufacturer curve selection for the adopted load assumptions.}
\begingroup\setlength{\tabcolsep}{3bp}
\begin{reporttable}{L{0.285714}L{0.285714}L{0.428571}}
\rowcolor{tablegrey} \textbf{Condition} & \textbf{Sensible load} & \textbf{Published curves} \\\hline
Maximum flow & 170.10 W/m2 & 40 and 60 Btu/(h ft2) \\
40\% flow & 68.04 W/m2 & 20 and 40 Btu/(h ft2) \\
\end{reporttable}\endgroup
The paired-curve method compares each directional ratio with both bounding load curves. A point beyond the published domain has no supported reading. Sections 3 and 4 use the louvered-diffuser selection interval. Confirming an ADPI value from the manufacturer curves would require a supported reading at each operating point.

\subhead{E.2 Linear-slot system}
The separate manufacturer-based system uses 63 Krueger Model 1900 linear-slot outlets. Its selection uses the Type 4 slot curves and the throw data for the selected slot configuration. The characteristic length runs in plan from the slot to the boundary or the midplane of the opposing outlet. Each outlet is checked at maximum and minimum flow.

The two systems use different outlets, layouts and throw curves. Drawings and schedules in Appendices A through D apply to the 80-outlet system. Applying the slot selection at 24/13 °C requires outlet-specific flows and a fresh curve check; the existing slot study does not establish performance at that operating point.

\tablenote{Manufacturer source: Krueger, 1900 Supply Published Performance Data, pp. 1, 4, 7 and 10. Air-distribution curves: reference [4].}
